\subsection{FreeRTOS}\label{sub:FreeRTOS} % (fold)
FreeRTOS is a class of Real Time Operating Systems(RTOS), that is designed to
be small enough for microcontrollers while not limiting its application to
purely microcontrollers. Out of the box, FreeRTOS provides core real time
scheduling functionality, inter-task communication, timing and synchronisation
primitives along with other possible add on functionalities. \cite{WhatFreeRTOSFreeRTOS2026}

FreeRTOS-MPU is a memory protected version of FreeRTOS, which as of writing is
supported on select ARM based microcontrollers. It enables tasks to run in
either privileged or unprivileged mode and restricts access to resources such
as RAM, executable code, peripherals and memory beyond the limit of a task's
stack. \cite{MemoryProtectionUnit2026}

Currently there exists a third party port of this implementation to the RISC-V
ISA \cite{FreeRTOSCommunitySupportedDemosRISCV_RV32_MPU_QEMU_VIRT_GCCmain}.
This is what forms the basis for this thesis implementation. There are a couple
of reason for this choice. The first and most important one is familiarity.
Previously I have worked with the RISC-V ISA and as such this reduces the
overhead needed on my end in regards to working with a different ISA. Second,
the related work on timing channels generally focuses on this ISA, so future
comparison between implementations are more relevant if all are built on the
same platform. At last, the \texttt{fence.t} instruction introduced in
Section~\ref{sub:temporal} is implemented with the RISC-V instruction set in
mind. While the underlying idea could be implemented on other platforms work on
hardware support is not as far along, so assuming this implementation will run
on real hardware at some point in the future, then RISC-V seems like the most
reasonable choice.

% subsection FreeRTOS (end)
